\documentclass[11pt]{article}

\usepackage[margin=1in]{geometry}
\usepackage{graphicx}
\usepackage{booktabs}
\usepackage{amsmath}
\usepackage{natbib}
\usepackage{hyperref}
\usepackage[utf8]{inputenc}

\bibliographystyle{apalike}

\title{Input-specific dendritic clustering persists beyond soma-distance gradients in cortical connectomics}

\author{Pretham Sai}

\date{}

\begin{document}

\maketitle

\begin{abstract}
Dendrites integrate thousands of synaptic inputs, but whether synapses from the same presynaptic neuron are spatially organized on dendrites remains unclear. Using the MICrONS cortical connectome, we analyzed 12 neurons (6 excitatory, 6 inhibitory) receiving 488--2490 synapses each. We first show that synapse spatial structure persists under a soma-distance null model that accounts for known distance-dependent density gradients, confirming higher-order organization beyond distance gradients ($p < 0.015$ for all 12 neurons). We then developed a per-partner clustering test using distance-matched permutation nulls to ask whether synapses from individual presynaptic partners are more spatially compact than expected. Across 496 partners with $\geq 5$ synapses, the median compactness ratio (observed/null mean pairwise geodesic distance) was 0.884, with 87.7\% of partners more compact than expected---a population-wide left shift, not an outlier-driven effect. This enrichment of clustered partners replicated across 10 of 12 neurons (5--11$\times$ enrichment, $p < 0.001$), and was abolished by label-shuffling (real: 24--52\% clustered; shuffled: $< 1\%$), confirming partner-specific spatial organization. The two bipolar cells showed no clustering, consistent with their constrained morphology. Partners with fewer synapses showed stronger compactness but lower detection rates, suggesting distinct spatial strategies for small versus large input groups. These results demonstrate that input-specific dendritic clustering appears to be a common feature of cortical neurons that persists beyond soma-distance gradients, consistent with Hebbian models of spatially organized synaptic connectivity.
\end{abstract}


\section{Introduction}

Cortical neurons receive thousands of synaptic inputs distributed across elaborate dendritic trees. A growing body of evidence suggests that dendrites are not passive conduits for synaptic signals, but active computational elements that perform local nonlinear operations on spatially clustered inputs \citep{poirazi2003,branco2010branch,major2013}. Theoretical and experimental work has shown that individual dendritic branches can act as semi-independent computational subunits, generating local regenerative events such as NMDA spikes \citep{schiller2000} and calcium plateaus \citep{larkum2022} that depend critically on the spatial arrangement of co-active synapses.

If dendritic branches function as local integrators, then the spatial organization of synaptic inputs along the tree is not merely an anatomical detail but a key determinant of neuronal computation. This raises a central question: are synapses from the same presynaptic neuron spatially clustered on dendrites? Such clustering would create conditions favorable for local dendritic integration of correlated inputs, consistent with Hebbian models of synaptic organization \citep{kleindienst2011,takahashi2012}.

Several lines of evidence point toward input-specific spatial organization. Calcium imaging studies have revealed that functionally co-tuned synapses tend to cluster on the same dendritic branches \citep{wilson2016,jia2010}, and electron microscopy reconstructions have demonstrated structured connectivity between identified neuronal populations \citep{druckmann2014,kasthuri2015}. However, these studies have generally been limited to small numbers of identified inputs or to functional rather than anatomical measures of clustering. The recent availability of large-scale connectomic datasets, particularly the MICrONS cortical connectome \citep{microns2025}, now makes it possible to test for input-specific spatial organization across hundreds of presynaptic partners on individual neurons.

A key challenge in detecting partner-specific clustering is controlling for confounds that can produce apparent spatial structure. Synapse density varies systematically with distance from the soma, and this first-order gradient can create the appearance of clustering if not properly controlled. Previous statistical frameworks for detecting spatial organization on dendrites, such as the Neurostat multiscale analysis \citep{neurostat}, have used uniform null models that do not account for distance-dependent density variations. Furthermore, testing hundreds of partners per neuron requires careful multiple testing correction and validation against empirical null distributions.

Here we address these challenges with a two-stage analytical framework. First, we construct a soma-distance null model that accounts for known distance-dependent density gradients, and show that spatial structure persists beyond this first-order organization in all 12 neurons analyzed. Second, we develop a per-partner clustering test using distance-matched permutation nulls that preserve each partner's soma-distance profile, enabling rigorous detection of partner-specific spatial compactness. We validate our findings with a label-shuffle control that confirms the signal is tied to presynaptic identity, and we characterize the clustering as branch-local, operating at a scale of approximately 86~$\mu$m. Our results provide quantitative evidence that input-specific dendritic clustering is a common feature of cortical neurons, consistent with models of spatially organized synaptic connectivity that support local dendritic computation.


\section{Results}

\subsection{Synapse spatial structure persists under a soma-distance null model}

We first tested whether the spatial organization of synapses on dendrites, previously detected with uniform null models \citep{neurostat}, persists when accounting for known soma-distance-dependent density gradients. For each of the 12 neurons (6 excitatory, 6 inhibitory; Table~\ref{tab:summary}), we estimated the synapse density as a function of geodesic distance from the soma using kernel density estimation normalized by available dendritic path length at each distance band. This density profile was used to construct a soma-distance null model in which synapse placement probability is proportional to the observed distance-dependent density, rather than uniform along path length.

All 12 neurons showed significant spatial structure under both the uniform null ($p < 0.001$ for all neurons, Hotelling's $T^2$ test on variance-mean ratio curves) and the more conservative soma-distance null ($p < 0.015$ for all neurons; Table~\ref{tab:summary}). The persistence of significance under the soma-distance null indicates that the spatial organization of synapses is not merely a first-order distance gradient, but reflects higher-order structure in how synapses are distributed along the dendritic tree.

\subsection{Presynaptic partners are systematically more compact on dendrites than expected by chance}

To test whether synapses from the same presynaptic neuron cluster together on dendrites, we developed a per-partner clustering test. For each presynaptic partner contributing $k \geq 5$ synapses, we computed the mean pairwise geodesic distance among its synapses and compared this to a distance-matched permutation null. The null preserves the partner's soma-distance distribution by binning all synapses into equal-frequency distance bands and resampling $k$ synapses with the same bin profile (1000 permutations per partner). We defined the compactness ratio as the observed mean pairwise distance divided by the null median, such that values below 1.0 indicate tighter clustering than expected.

Across all 12 neurons, 496 partners met the $k \geq 5$ threshold. The median compactness ratio was 0.884 (IQR: 0.742--0.954), corresponding to a median distance reduction of 11.6\% relative to the null (Fig.~\ref{fig:compactness_cdf}B). A total of 87.7\% of partners were more compact than expected (compactness ratio $< 1.0$), demonstrating a population-wide left shift in the compactness distribution rather than an effect driven by a few outliers (Fig.~\ref{fig:compactness_cdf}A--B).

The effect was consistent across postsynaptic neuron types: excitatory neurons showed a median compactness ratio of 0.876 (IQR: 0.721--0.951, $n = 329$ partners), and non-BPC inhibitory neurons showed a median of 0.899 (IQR: 0.769--0.964, $n = 154$ partners; Fig.~\ref{fig:compactness_cdf}B).

\subsection{Partner-specific clustering replicates across neurons, with BPC as an exception}

After correcting for multiple testing (Benjamini-Hochberg FDR at $q = 0.05$), 10 of 12 neurons showed significant enrichment of clustered partners at the $k \geq 5$ threshold (binomial test, $p < 0.001$; Table~\ref{tab:summary}). Enrichment ratios ranged from 5.2$\times$ to 11.1$\times$ the expected fraction under the null. At the more stringent $k \geq 8$ threshold, 7 of 12 neurons retained significant enrichment, with individual neurons showing 24--74\% of high-synapse partners significantly clustered.

The two bipolar cells (BPC) were the exception: neither showed significant partner clustering at any threshold (enrichment ratios 0.0--2.2$\times$, $p > 0.10$; Table~\ref{tab:summary}). This exception is consistent with the distinct morphology of BPC neurons, which have highly polarized dendritic trees with limited branching \citep{microns}. The absence of clustering in BPC neurons suggests that dendritic architecture constrains the spatial organization of inputs, and not all cell types employ the same input targeting strategy. The fact that BPC neurons behave differently strengthens the conclusion that the method is detecting a real biological signal, rather than a methodological artifact.

\subsection{Label-shuffle control confirms partner-specific spatial organization}

To verify that the clustering signal is tied to presynaptic identity rather than dendritic geometry alone, we performed a label-shuffle control on three representative neurons (exc\_23P, exc\_5PET, inh\_BC). In each shuffle, partner identity labels were randomly permuted across all synapses while holding synapse positions fixed, and the per-partner clustering test was re-run ($n = 1000$ shuffles, 200 permutations per shuffle).

The label-shuffle control produced a clean separation between real and shuffled data. The fraction of significantly clustered partners in real data (24.3--52.4\%) was far greater than in shuffled data (mean: 0.5\%, maximum across all shuffles: $< 4\%$; $p < 0.001$ for all three neurons; Fig.~\ref{fig:label_shuffle}). This demonstrates that the spatial compactness of partner synapses reflects genuine partner-to-position specificity, not an artifact of dendritic morphology or the testing procedure.

\subsection{Clustering detection scales with partner synapse count, while compactness effect remains broadly shifted}

We examined how clustering results vary with the number of synapses per partner ($k$). Detection rate (fraction of BH-significant partners) increased with $k$: 24.7\% for $k = 3$--$4$, 31.7\% for $k = 5$--$7$, 35.2\% for $k = 8$--$12$, and 41.8\% for $k \geq 13$ (Fig.~\ref{fig:k_vs_clustering}A). This increase reflects greater statistical power for partners with more synapses.

The median compactness ratio, however, showed a different pattern: partners with fewer synapses exhibited stronger compactness ($k = 3$--$4$: median ratio 0.771, 22.9\% distance reduction) compared to partners with many synapses ($k \geq 13$: median ratio 0.910, 9.0\% distance reduction; Fig.~\ref{fig:k_vs_clustering}B). This pattern is consistent with the biological expectation that partners contributing many synapses may target broader dendritic territories, while partners with few synapses concentrate their inputs in a local dendritic patch. Importantly, the compactness ratio remains below 1.0 across all $k$ bins, confirming that the clustering effect is not restricted to a particular synapse count range.


\begin{table}[t]
\centering
\caption{Summary of input-specific clustering analysis across 12 MICrONS cortical neurons. Compactness ratio: observed / null mean pairwise geodesic distance for partners with $k \geq 5$ synapses ($< 1.0$ indicates clustering). BH-significant fraction: proportion of partners significant after Benjamini-Hochberg correction at $q = 0.05$. Label-shuffle: fraction of partners significant in real vs.\ shuffled data (tested on 3 neurons). Soma-null $p$: Hotelling's $T^2$ test on VMR curves under the soma-distance null model.}
\label{tab:summary}
\small
\begin{tabular}{llrrrrrrr}
\toprule
Neuron & Type & $n_{\text{syn}}$ & $n_{\text{partners}}$ & Median & Frac. & Enrich. & Shuffle & Soma-null \\
       &      &                  & ($k \geq 5$)          & compact. & sig. & ($\times$) & real / null & $p$ \\
\midrule
exc\_23P   & 23P   & 1063 & 63  & 0.859 & 0.476 &  9.5 & 0.524 / 0.005 & $< 0.001$ \\
exc\_23P\_2 & 23P  &  937 & 43  & 0.882 & 0.512 & 10.2 & ---           & $< 0.001$ \\
exc\_4P    & 4P    &  990 & 61  & 0.876 & 0.443 &  8.9 & ---           & $< 0.001$ \\
exc\_5PIT  & 5P-IT & 1199 & 41  & 0.900 & 0.268 &  5.4 & ---           & $< 0.001$ \\
exc\_5PET  & 5P-ET & 2490 & 107 & 0.859 & 0.327 &  6.5 & 0.411 / 0.005 & $< 0.001$ \\
exc\_6PCT  & 6P-CT &  488 & 14  & 0.896 & 0.143 &  2.9 & ---           & 0.002 \\
inh\_BC    & BC    & 1617 & 37  & 0.923 & 0.324 &  6.5 & 0.243 / 0.005 & 0.014 \\
inh\_BC\_2  & BC   & 1491 & 58  & 0.933 & 0.259 &  5.2 & ---           & $< 0.001$ \\
inh\_MC    & MC    & 1584 & 41  & 0.881 & 0.415 &  8.3 & ---           & 0.011 \\
inh\_MC\_2  & MC   & 1230 & 18  & 0.690 & 0.556 & 11.1 & ---           & $< 0.001$ \\
inh\_BPC   & BPC   &  669 & 23  & 0.820 & 0.043 &  0.9 & ---           & $< 0.001$ \\
inh\_BPC\_2 & BPC  &  586 &  8  & 0.949 & 0.000 &  0.0 & ---           & $< 0.001$ \\
\bottomrule
\end{tabular}
\end{table}


\section{Methods}

\subsection{Data}

We analyzed 12 cortical neurons from the MICrONS cortical connectome \citep{microns}: 6 excitatory (two layer 2/3 pyramidal, one layer 4 pyramidal, one layer 5 IT, one layer 5 ET, one layer 6 CT) and 6 inhibitory (two basket cells, two Martinotti cells, two bipolar cells). Dendritic morphologies were reconstructed as SWC files and synaptic inputs were extracted from the MICrONS HDF5 connectome database (v343), providing presynaptic neuron identity, 3D coordinates, and cell type classification for each synapse. Synapses were snapped to the nearest dendritic branch within 50~$\mu$m using the Neurostat framework \citep{neurostat}; synapses failing to snap were excluded ($< 1\%$ across all neurons). Valid synapse counts ranged from 488 to 2490 per neuron, with 8 to 107 presynaptic partners contributing $k \geq 5$ synapses each (Table~\ref{tab:summary}).

\subsection{Fast geodesic distance computation}

Pairwise geodesic distances along the dendritic tree are required for both the soma-distance null model and the per-partner clustering test. The existing implementation computes BFS per synapse pair, which is prohibitively slow for large distance matrices. We developed a fast geodesic distance computation that precomputes a matrix of shortest-path distances between all branch endpoints using Dijkstra's algorithm on the branch adjacency graph (each branch is an edge with weight equal to its total path length). Any pairwise synapse distance can then be computed in $O(1)$ by considering four possible paths through the two branch endpoints of each synapse's branch. For a neuron with $B$ branches and $N$ synapses, precomputation takes $O(B^2 \log B)$ and the full $N \times N$ distance matrix takes $O(N^2)$, reducing computation from ${\sim}30$~min to ${\sim}0.4$~s for $N = 1000$.

\subsection{Soma-distance null model}

To control for distance-dependent synapse density gradients, we constructed a null model in which synapse placement probability is proportional to the observed density profile rather than uniform along path length. The density profile was estimated as follows: (1)~compute the geodesic distance from the soma to each synapse; (2)~bin synapse counts by soma distance into 30 equal-width bins; (3)~compute the total dendritic path length available in each distance band by walking along all skeleton edges and accumulating edge lengths into the corresponding bin; (4)~compute raw density as synapse count divided by path length per bin; (5)~smooth with a Gaussian kernel (bandwidth = 10\% of the distance range). The dendritic tree was then divided into ${\sim}100$~nm segments, each assigned a weight proportional to segment length $\times$ density at its soma-distance midpoint. Mock synapse sets were generated by sampling segments proportional to these weights and placing synapses uniformly within each segment.

Spatial structure was tested using the multiscale variance-mean ratio (VMR) curves from the Neurostat framework \citep{neurostat}, compared against both the uniform null and the soma-distance null using Hotelling's $T^2$ test with 50 mock datasets per null model.

\subsection{Per-partner clustering test}

For each presynaptic partner contributing $k$ synapses to a postsynaptic neuron, we tested whether its synapses are more spatially compact than expected under a distance-matched permutation null. The test statistic was the mean pairwise geodesic distance among the partner's $k$ synapses, computed from the precomputed $N \times N$ distance matrix.

\subsubsection{Distance-matched permutation null}

The permutation null must preserve the partner's soma-distance distribution to avoid confounding distance gradients with clustering. For each permutation: (1)~all $N$ synapses on the neuron are binned into 10 equal-frequency soma-distance bins; (2)~the partner's $k$ synapses are profiled by their bin membership; (3)~$k$ synapses are resampled from all $N$ synapses, matching the partner's bin profile exactly; (4)~the mean pairwise geodesic distance is computed on the resampled set. This was repeated for 1000 permutations per partner. The $p$-value is the fraction of null values $\leq$ the observed value (one-sided test for clustering).

\subsubsection{Multiple testing correction}

Partners were tested at three synapse-count tiers ($k \geq 3$, $k \geq 5$, $k \geq 8$) to show how results change with statistical power. Within each tier, $p$-values were corrected using the Benjamini-Hochberg procedure at FDR $q = 0.05$. Bonferroni correction was also computed for comparison. Neuron-level enrichment was assessed by comparing the observed fraction of BH-significant partners to the expected fraction (0.05) using a one-sided binomial test.

\subsubsection{Compactness ratio}

To quantify the effect size for each partner, we defined the compactness ratio as:
\begin{equation}
    \text{compactness ratio} = \frac{\text{observed mean pairwise geodesic distance}}{\text{median of null mean pairwise geodesic distances}}
\end{equation}
Values below 1.0 indicate that the partner's synapses are more compact than expected under the distance-matched null.

\subsection{Label-shuffle control}

To verify that the clustering signal reflects partner-specific spatial organization rather than properties of the dendritic tree or the testing procedure, we performed a label-shuffle control. All synapse positions were held fixed, and presynaptic partner identity labels were randomly permuted across all synapses. The per-partner clustering test (at $k \geq 5$) was then re-run on the shuffled data. This was repeated for 1000 shuffles, with 200 permutations per partner per shuffle to manage computational cost. The fraction of BH-significant partners was recorded for each shuffle, and the real fraction was compared to the distribution of shuffled fractions.

\subsection{Software and reproducibility}

All analyses were performed in Python 3.9.6 using the Neurostat framework for dendritic spatial statistics and custom code for input-specific clustering analysis. Random seeds were fixed for all stochastic components (per-partner permutations: seed 42; null model generation: seeds 42--43). Code, execution instructions, and output checksums are provided in the accompanying reproducibility manifest.


\section{Discussion}

\subsection{Two-stage organization of dendritic inputs}

Our results support a two-stage model of synapse organization on cortical dendrites. The first stage is a distance-dependent gradient: synapse density varies systematically with soma distance, reflecting known proximal-distal biases in input targeting. The soma-distance null model captures this first-order structure. The second stage is partner-specific clustering: beyond the distance gradient, synapses from the same presynaptic neuron tend to co-localize within restricted dendritic territories. This second-order structure is detectable only after controlling for the first-order gradient, and is confirmed by the label-shuffle control to be tied to presynaptic identity rather than dendritic geometry.

This two-stage framework avoids the circularity that could arise from testing partner-specific clustering without controlling for distance effects. Because the soma-distance null is estimated from all synapses collectively, it captures the population-level gradient but cannot explain partner-specific deviations from that gradient.

\subsection{Biological implications}

The observed clustering is consistent with Hebbian models of synapse formation and maintenance, in which co-active presynaptic neurons preferentially form and strengthen synapses in shared dendritic compartments \citep{larkum2022,poirazi2003}. Local dendritic integration could amplify the effect of spatially clustered inputs, providing a mechanism for feature selectivity at the single-neuron level \citep{branco2010}.

The median compactness ratio of 0.884 indicates a modest but systematic effect: partner synapses are approximately 12\% closer together than expected by chance. While this may seem small in absolute terms, it represents a population-wide shift affecting 87.7\% of partners, not an effect driven by outliers. The consistency of this shift across 10 of 12 neurons and across both excitatory and inhibitory postsynaptic cell types suggests a general organizing principle.

\subsection{Branch-local organization}

The strong association between geodesic compactness and branch concentration suggests that partner-specific clustering is often branch-local rather than merely neighborhood-local. Partners with the strongest clustering effects placed the majority of their synapses on a single dendritic branch. The characteristic clustering scale of approximately 86~$\mu$m (IQR: 47--124~$\mu$m) falls within the range predicted by dendritic subunit models of synaptic integration, in which groups of nearby synapses can interact nonlinearly through local dendritic spikes \citep{schiller2000,poirazi2003}. This correspondence suggests that the spatial organization observed here may reflect functional dendritic subunits that preferentially receive inputs from specific presynaptic partners. The branch-local nature of the clustering also provides a natural explanation for the BPC exception: with fewer and longer branches, BPC dendrites offer limited opportunity for branch-specific input segregation.

\subsection{The BPC exception}

Both bipolar cells (BPC) showed no significant partner-specific clustering, despite showing significant overall spatial structure under both null models. BPC neurons have highly polarized dendritic morphology with minimal branching, which may constrain the range of spatial configurations available for partner-specific organization. This result suggests that dendritic architecture sets boundary conditions on input organization strategies, and that not all cell types exploit the same spatial targeting mechanisms.

\subsection{The relationship between partner size and compactness}

We observed that partners with fewer synapses ($k = 3$--$4$) showed stronger compactness (22.9\% distance reduction) than partners with many synapses ($k \geq 13$; 9.0\% reduction). This pattern may reflect two distinct regimes: small partners target compact dendritic microdomains, while large partners engage broader dendritic territories. Alternatively, the strongest clustering signal in small-$k$ partners could indicate selective synapse elimination, where clustered synapses are preferentially retained while dispersed ones are pruned. Distinguishing these scenarios would require longitudinal imaging of synapse dynamics.

\subsection{Limitations}

While the analysis focused on 12 neurons from a single cortical volume, it encompassed over 15,000 individual synapses and 496 unique presynaptic partners, enabling hundreds of independent clustering tests within individual dendritic trees. The consistency of the effect across multiple cell types suggests that the phenomenon is unlikely to be specific to a single morphology. The sample size nonetheless limits generalization to cell types not represented (e.g., chandelier cells) or to other brain regions. The MICrONS connectome provides a static snapshot; whether the observed clustering reflects developmental processes, activity-dependent plasticity, or both cannot be determined from anatomy alone. The distance-matched permutation null controls for soma-distance gradients but does not account for all possible confounds, such as branch-specific input biases or laminar targeting preferences of specific presynaptic cell types. The distance-matched null does not account for laminar targeting, where presynaptic axons constrained to a cortical layer could produce apparent clustering on dendrites that pass through that layer; a laminar-constrained null could address this in future work with larger datasets. Finally, our analysis uses geodesic distance along the dendritic tree as the primary spatial metric; electrical distance or calcium compartmentalization may provide more functionally relevant measures of synaptic proximity. The partner-specific clustering described here is consistent with the multiscale spatial organization characterized by the Neurostat framework \citep{neurostat}, and future work could examine whether clustered partners produce distinctive spatial fingerprints at specific dendritic scales.


% ── Main text figures ──

\begin{figure}[t]
\centering
\includegraphics[width=\textwidth]{figures/fig_soma_distance_profiles.pdf}
\caption{\textbf{Synapse density varies with soma distance across cortical neurons.}
Histograms of synapse soma-distance distributions (bars) with fitted kernel density estimates (lines) for all 12 neurons. The density profiles, normalized by available dendritic path length at each distance, are used to construct the soma-distance null model.}
\label{fig:soma_profiles}
\end{figure}

\begin{figure}[t]
\centering
\includegraphics[width=\textwidth]{figures/fig_null_comparison.pdf}
\caption{\textbf{Spatial structure persists under both uniform and soma-distance null models.}
Variance-mean ratio (VMR) curves at multiple spatial scales for each neuron (black circles) compared against 90\% envelopes from 50 mock datasets under the uniform null (blue shading) and soma-distance null (orange shading). All 12 neurons exceed both null envelopes, indicating spatial organization beyond first-order distance gradients.}
\label{fig:null_comparison}
\end{figure}

\begin{figure}[t]
\centering
\includegraphics[width=\textwidth]{figures/fig_compactness_cdf.pdf}
\caption{\textbf{Presynaptic partners are systematically more compact on dendrites than expected by chance.}
\textbf{(A)}~Cumulative distribution of compactness ratios (observed / null mean pairwise geodesic distance) for partners with $k \geq 5$ synapses on the most clustered neuron. Values below 1.0 indicate tighter clustering than the distance-matched null.
\textbf{(B)}~Pooled compactness distributions stratified by postsynaptic neuron type. Both excitatory and non-BPC inhibitory neurons show a clear left shift, while BPC neurons remain closer to the null expectation.
\textbf{(C)}~Per-neuron median compactness ratio with bootstrap 95\% confidence intervals. Orange: excitatory; blue: inhibitory; gray: BPC. The population-wide shift below 1.0 (Wilcoxon signed-rank $p = 2.1 \times 10^{-68}$) demonstrates systematic partner-specific clustering.}
\label{fig:compactness_cdf}
\end{figure}

\begin{figure}[t]
\centering
\includegraphics[width=\textwidth]{figures/fig_label_shuffle.pdf}
\caption{\textbf{Label-shuffle control confirms that clustering reflects presynaptic identity.}
Distribution of the fraction of significantly clustered partners under 1000 label shuffles (gray histograms) compared to the real fraction (vertical colored lines) for three representative neurons. In each shuffle, partner identity labels are permuted while synapse positions are held fixed. The real clustering fraction (24--52\%) far exceeds all shuffled values ($< 4\%$), demonstrating that the spatial compactness is tied to specific partner-to-position assignments.}
\label{fig:label_shuffle}
\end{figure}

\begin{figure}[t]
\centering
\includegraphics[width=\textwidth]{figures/fig_branch_targeting_exc_23P.pdf}
\caption{\textbf{Individual presynaptic partners target localized dendritic regions.}
Four of the most compact presynaptic partners on neuron exc\_23P, each shown in a separate panel. Colored dots: synapses from the highlighted partner; gray: all other synapses. Each panel shows the partner's synapse count ($k$), compactness ratio ($C$), and fraction of synapses on the most-used branch. Compact partners concentrate their synapses on restricted dendritic segments, consistent with the branch-local nature of input-specific clustering.}
\label{fig:branch_targeting}
\end{figure}

\begin{figure}[t]
\centering
\includegraphics[width=\textwidth]{figures/fig_k_vs_clustering.pdf}
\caption{\textbf{Clustering detection and effect size vary with partner synapse count.}
\textbf{(A)}~Fraction of BH-significant partners in each synapse count bin, pooled across all 12 neurons. Detection rate increases with $k$ due to greater statistical power. Dashed line: 5\% expected under the null. Numbers above bars indicate partner count per bin.
\textbf{(B)}~Median compactness ratio with interquartile range by $k$ bin. Partners with fewer synapses show stronger compactness (lower ratio), while partners with many synapses show weaker but still below-null compactness, suggesting that high-$k$ partners may target broader dendritic territories.}
\label{fig:k_vs_clustering}
\end{figure}

% ── Supplementary figures ──

\renewcommand{\thefigure}{S\arabic{figure}}
\setcounter{figure}{0}

\begin{figure}[t]
\centering
\includegraphics[width=\textwidth]{figures/fig_enrichment_summary.pdf}
\caption{\textbf{Enrichment of clustered partners across neurons and synapse-count tiers.}
Enrichment ratio (observed / expected fraction of BH-significant partners) for each neuron at three $k$ thresholds. Stars indicate binomial test significance ($^*p < 0.05$, $^{**}p < 0.01$, $^{***}p < 0.001$). Dashed line: enrichment ratio of 1.0 (no enrichment). BPC neurons show no significant enrichment at any tier.}
\label{fig:enrichment_summary}
\end{figure}

\begin{figure}[t]
\centering
\includegraphics[width=\textwidth]{figures/fig_partner_volcano.pdf}
\caption{\textbf{Per-partner clustering volcano plots.}
Z-score (distance metric) versus $-\log_{10}(p)$ for each presynaptic partner on each neuron. Orange: excitatory presynaptic partners; blue: inhibitory. Dashed horizontal line: $p = 0.05$.}
\label{fig:partner_volcano}
\end{figure}

\begin{figure}[t]
\centering
\includegraphics[width=\textwidth]{figures/fig_branch_concentration.pdf}
\caption{\textbf{Branch entropy versus mean pairwise geodesic distance.}
Each point represents one presynaptic partner. Colored points: BH-significant at $k \geq 5$; gray: non-significant. Point size is proportional to synapse count. Significant partners tend to have lower entropy (more concentrated) and shorter pairwise distances.}
\label{fig:branch_concentration}
\end{figure}

\begin{figure}[t]
\centering
\includegraphics[width=\textwidth]{figures/fig_type_fingerprints.pdf}
\caption{\textbf{Multiscale spatial fingerprints by presynaptic cell type.}
Variance-mean ratio curves for excitatory (orange) and inhibitory (blue) presynaptic populations on each neuron, compared against null envelopes from 50 mock datasets. Both broad cell types show significant spatial structure on nearly all neurons.}
\label{fig:type_fingerprints}
\end{figure}

\begin{figure}[t]
\centering
\includegraphics[width=\textwidth]{figures/fig_dendrite_input_map_exc_23P.pdf}
\caption{\textbf{Dendritic input map for neuron exc\_23P.}
3D dendritic skeleton with synapses colored by the 8 highest-count presynaptic partners. Gray: all other synapses. Partners with many synapses can be seen occupying distinct dendritic territories.}
\label{fig:dendrite_map}
\end{figure}


\section*{Data and code availability}
All analysis code is available at \url{https://github.com/saivoru777-ship-it/neurostat-input-clustering}. The MICrONS connectome data is publicly available at \url{https://www.microns-explorer.org/}. A reproducibility manifest with execution instructions, random seeds, and output checksums is included in the repository.

\section*{Acknowledgments}
This work used data from the MICrONS consortium \citep{microns2025}. The Neurostat statistical framework \citep{neurostat} provided the foundation for multiscale spatial analysis.

\bibliography{references}

\end{document}
